\chapter{Introduction}

\begin{chapterintro}
This chapter introduces unemployment and the methodology used to approach this research.
\end{chapterintro}

Unemployment is a worldwide crisis affecting hundreds and thousands, with the global unemployment rate at 5.7\% over 2015--2016 and at an alarming 26.7\% in South Africa in 2016. The United Nations thus rightly listed “Decent Work and Economic Growth” as goal number eight of the Sustainable Development Goals (SDGs) with the goal to achieve full and productive employment for both men and women through job diversification and technological upgrade and innovation. Unemployment is thus categorised as an Information Communication Technology for Development (ICT4D) problem; characterized by issues of low-income, illiteracy, marginalisation and other forms of underdevelopment. As expected, unemployment results in several social ills such as increased crime, ill-health and educational inequalities thus widening the gap between the rich and the poor. In addressing unemployment, several applications have been created to tackle the issue, these can be classified as ICT4D interventions. Among these are Ummeli, a website for finding jobs for the unemployed in South Africa and Babajob, a website where job seekers are connected to employers in India.

Job readiness (JR defined in Section 2.2) programs and non-governmental organisations (NGOs) in Cape Town address the unemployment problem through upskilling but lack ICT4D interventions. JR graduates still face issues such as poor search skills, lack of mobile devices and data to access job opportunities. Thus, Co-design features as an approach to offers an opportunity to work with these NGOs to identify where ICT4D expertise might play a role and offers a platform to work directly with the affected stakeholders to design and evaluate solutions. In this research, we used Co-design to understand unemployment and explore ways in which ICTs can be used to ease job searching for job seekers. By using Co-design, we are able to make use of a combination of different tools and methods to engage our users, flexibly move between cycle phases and gain rich user-oriented suggestions and solutions. Thus, we interchanged between the use of participant observation, surveys, one-on-one discussions and workshops as was needed for the different enquiry and design stages.

The first cycle of our research was focused on discovering and understanding our users’ capabilities and their usage of communication and networking technologies, which are used in job searching and other everyday tasks. In Cycle 2, we worked together with our users to (1) design an enabling artefact to assist with job finding, and then (2) develop a working version of the artefact. Finally, in Cycle 3, we deployed the improved version of the artefact for a period of two weeks and assessed its feasibility and potential impact.

Through our research, we worked with job seekers, who we recruited at two JR training sites as well as their trainers to grasp the issues around unemployment in Cape Town, South Africa. From our findings, we discuss the combination of technologies and approaches to tackle the problem; mobile and web-based methods, door-to-door and electronic searches, SMS and email communications for job searching for the low-skilled. We discovered our users’ preferences as well as the barriers to using the preferred mode of communication. Finally, we deployed an application which combined some of our users’ different means of communicating and job finding, as well as existing best practices of job search sites.

Our two NGO partner organisations: Afrika Tikkun and The Zanokhanyo Network are both organisations focused on attracting unemployed individuals who are seeking to become more employable, by providing them with JR classes, which are classes which prepare them for the world of works. Our research involved two different organisations to improve the rigour of our outcomes. Our participants were: job-seeking students as beneficiaries, their trainers as intermediaries and the researcher as an observer, interviewer, data analyst and software developer. All three stakeholders worked together as Co-designers in the design process of the eventual artefact.

We discovered what worked and what did not work for the low-skilled job seekers and went on to develop and deploy an application to access its practical feasibility. Through the research, several trade-offs affected our findings and thus we list them for guidelines for future work with unemployed individuals and future work in Co-design.

\section{Research Question}

{\setlength{\parindent}{0pt}
\setlength{\parskip}{1em}
\setlist[enumerate]{leftmargin=*, font=\bfseries}

In this research, we sought to understand the ICT setbacks for employment amid the low-skilled. Thus, we ask:

\begin{enumerate}[label=RQ\arabic*:]
    \item \textit{In what ways are the low-skilled using existing ICTs to find jobs?}
\end{enumerate}

Question one allows us to understand how existing ICTs fit into the lives of job seekers so as not to bring in off-the-shelf solutions which are not usable for our users.

\begin{enumerate}[label=RQ\arabic*:, resume]
    \item \textit{What are the barriers to the use of ICTs for job searching for the low-skilled?}
    \item \textit{Which technology or technologies are appropriate for low-skilled job seekers to find jobs?}
\end{enumerate}

Question two and three then help us to identify the gaps in the current technologies and go on to identify new technologies or new combinations of existing technologies to bridge the gap. To bridge the barriers or gaps we interact with low-skilled job seekers to determine what is appropriate for them by discussing their experiences, preferences and abilities through a series of tools and techniques.

Then we ask:

\begin{enumerate}[label=RQ\arabic*:, resume]
    \item \textit{What are the salient design requirements for ICTs used to assist job seekers in job finding?}
\end{enumerate}

Through the various interactions we were able to determine which issues arose more frequently, i.e.\ which were salient, thus, we contributed to understanding job searching among the low-skilled and other marginalised users with similar technology limitations.

Finally, we asked:

\begin{enumerate}[label=RQ\arabic*:, resume]
    \item \textit{What are the trade-offs to be made when performing Co-design with low-skilled, unemployed or marginalised individuals from NGOs?}
\end{enumerate}

Because of the various methods used to interact with our users, we were able to discuss the trade-offs, in other words: which techniques or approaches worked better under which circumstances and to what extent various methods should or should not be applied. Thus, contributing to understanding Co-design and NGO engagement in ICT4D projects.
}